\section{Resumen}

\subsection{Resumen}

El presente proyecto propone el diseño y desarrollo de \textbf{SocioUnido}, una plataforma E2E (Empresa a Empresa) \textit{SaaS (Software como Servicio)} orientada a modernizar la gestión administrativa, optimizar la recaudación y asegurar el control de accesos en clubes deportivos argentinos. La solución aborda tres ineficiencias críticas: la morosidad societaria silenciosa, el fraude en los ingresos por préstamo de credenciales y el desaprovechamiento de activos institucionales.

El ecosistema integra un panel web administrativo y una aplicación web progresiva (PWA) para socios y trabajadores. Tecnológicamente, destaca un motor predictivo (\textit{Machine Learning}) para prevenir bajas societarias, un asistente conversacional vía WhatsApp con Procesamiento de Lenguaje Natural (PLN) para pagos, y proyecta a futuro un control de accesos (\textit{Smart Access}) mediante tokens TOTP con validación 100\% \textit{offline}.

El desarrollo aplica la metodología ágil Scrum, pruebas automatizadas y despliegue continuo (CI/CD). La estrategia comercial (\textit{Go-To-Market}), validada mediante las 5 C y 4 P, prioriza inicialmente a los clubes del ascenso metropolitano y del interior para consolidar la herramienta. En conclusión, SocioUnido promueve el saneamiento financiero, la inclusión tecnológica y la reducción del impacto ambiental al desmaterializar credenciales físicas.

\subsection{Palabras Clave}

Plataforma SaaS (Software como Servicio), Clubes de Fútbol, Control de Accesos Fuera de Línea, Aprendizaje Automático, Procesamiento de Lenguaje Natural (PLN), Algoritmo TOTP (Contraseña de un Solo Uso Basada en Tiempo), Metodología Scrum, Gestión Deportiva.

\newpage

\section{Abstract}

\subsection{Abstract}

This project presents the design and development proposal for \textbf{SocioUnido}, a Business-to-Business (B2B) Software as a Service (SaaS) platform engineered to modernize administrative management, optimize revenue collection, and secure access control in Argentine sports clubs. The solution addresses three critical inefficiencies: undetected member delinquency, gate fraud via credential sharing, and the underutilization of institutional assets.

The ecosystem integrates an administrative web dashboard and a progressive web application (PWA) for members. Technologically, it highlights a predictive engine (Machine Learning) to prevent member churn, a conversational WhatsApp assistant with Natural Language Processing (NLP) for payments, and projects a future offline access control system (\textit{Smart Access}) using TOTP tokens.

The development applies the agile Scrum methodology, automated testing, and continuous deployment (CI/CD). The commercial Go-To-Market strategy, validated through the 5 Cs and 4 Ps, initially prioritizes lower-tier clubs to solidify the tool. In conclusion, SocioUnido promotes financial recovery, technological inclusion, and environmental impact reduction by dematerializing physical credentials.

\subsection{Keywords}

SaaS Platform, Football Clubs, Offline Access Control, Machine Learning, Natural Language Processing (NLP), TOTP Algorithm, Scrum Methodology, Sports Management.
